% ============================================================
%  系统开发工具基础 · 实验报告 LaTeX 模板
%  作者：曾奕豪
%  编译方式：XeLaTeX（需 ctex 宏包支持中文）
%  用法：每周复制本文件为 weekN.tex，仅修改周次、日期、内容即可
%  说明：本模板不依赖任何图片即可编译；
%        图片文件存在时自动插入，不存在时显示占位提示框。
% ============================================================

\documentclass[12pt,a4paper]{article}

% ---------- 中文字体与页面 ----------
% 中文支持：使用 ctex 宏包（编译用 XeLaTeX）。
% 若在 Linux/TeX Live 下缺少 ctex，可安装 texlive-lang-chinese。
\usepackage[UTF8]{ctex}
\usepackage[margin=2.5cm]{geometry}
\usepackage{graphicx}
\graphicspath{{../figures/}}   % 图片统一放在 paper/figures/ 下
\usepackage{float}
\usepackage{booktabs}
\usepackage{xcolor}
\usepackage{listings}
\usepackage{enumitem}
\usepackage[colorlinks=true,linkcolor=black,urlcolor=blue]{hyperref}

% ---------- 页眉 / 页脚 ----------
\usepackage{fancyhdr}
\pagestyle{fancy}
\fancyhf{}
\fancyhead[R]{实验报告 · 第 \Week 周}
\fancyfoot[C]{\thepage}
\renewcommand{\headrulewidth}{0.4pt}

% ---------- 代码高亮环境 ----------
\lstset{
  basicstyle=\ttfamily\small,
  keywordstyle=\color{blue!70!black}\bfseries,
  commentstyle=\color{gray!70!black}\itshape,
  stringstyle=\color{red!60!black},
  numbers=left,
  numberstyle=\tiny\color{gray},
  frame=single,
  breaklines=true,
  tabsize=4,
  captionpos=b,
  showstringspaces=false,
}

% ---------- 图片自动插入宏 ----------
% 题内图片：图注自动按「图一、图二...」编号。
%   \resetpic              每题开头调用一次，重置编号
%   \resultimg[宽度]{图片文件名}   图片存在→插入；不存在→显示占位框（不报错）
% 自定义图注（如 commit 截图）：
%   \resultimgcap[宽度]{图片文件名}{图注}
\newcounter{picnum}
\newcommand{\resetpic}{\setcounter{picnum}{0}}
\newcommand{\resultimg}[3][0.8\textwidth]{%
  \IfFileExists{../figures/\detokenize{#2}}{%
    \begin{figure}[H]
      \centering
      \includegraphics[width=#1]{../figures/\detokenize{#2}}
      \parbox{0.8\textwidth}{\centering\small #3}
    \end{figure}%
  }{%
    \begin{center}
      \fbox{\parbox{0.8\textwidth}{\centering\itshape [图片待补充] \detokenize{#2}\\\small #3}}
    \end{center}%
  }%
}
\newcommand{\resultimgcap}[3][0.8\textwidth]{%
  \IfFileExists{../figures/\detokenize{#2}}{%
    \begin{figure}[H]
      \centering
      \includegraphics[width=#1]{../figures/\detokenize{#2}}
      \parbox{0.8\textwidth}{\centering\small #3}
    \end{figure}%
  }{%
    \begin{center}
      \fbox{\parbox{0.8\textwidth}{\centering\itshape [图片待补充] \detokenize{#2}\\\small #3}}
    \end{center}%
  }%
}

% ---------- 文档信息（每周只需改这处） ----------
\newcommand{\Week}{1}                        % 周次
\newcommand{\ReportTitle}{系统开发工具基础实验报告}
\newcommand{\ReportDate}{\today} % 日期，编译时自动取当天
\newcommand{\RepoLink}{https://github.com/hzydy00/Base-of-system-development-tools} % 仓库链接

% ============================================================
\begin{document}

% ---------- 封面 / 标题区 ----------
\begin{center}
  {\LARGE\bfseries \ReportTitle}\\[6pt]
  {\Large 第 \Week\ 周}\\[12pt]
  {\large 姓名：曾奕豪}\\
  {\large 课程：系统开发工具基础}\\[6pt]
  {\large \ReportDate}\\[16pt]
  {\large 仓库：\url{\RepoLink}}
\end{center}
\vspace{0.8em}
\hrule
\vspace{1.2em}

% ---------- 1. 课上内容与结果 ----------
\section{课上内容与结果}
% 开头先简述本周实验的主题、目标与涉及的工具 / 技术，再逐题列出课上内容。
% 课上完成 4 个不同主题的题目，每题按"题目要求 → 代码 → 运行结果"组织。
课程概览+Shell入门；VersionControl and Git；LaTeX 文档编辑

\subsection*{实验环境}
本实验基于 Windows 11 操作系统（版本 10.0.26200），使用 Git Bash 作为命令行环境，其内置 Git 版本为 2.53.0.windows.2。
实验涉及的 Shell 操作均在 Git Bash 中完成；版本控制操作使用 Git 命令行工具完成；LaTeX 文档采用 MiKTeX 26.5 发行版中的 XeLaTeX 编译，并通过 ctex 宏包处理中文排版。

\subsection*{题目一}
\resetpic
% 题目一 问题截图（代码汇总之前）
\resultimgcap[0.8\textwidth]{week1_q1.png}{题目一 题目要求}

\begin{lstlisting}[language=bash]
# TODO: 题目一代码汇总

# 1. 创建 input/docs 和 input/tmp 子目录
mkdir -p q01/input/docs q01/input/tmp

# 2. 创建带空格的文件 notes one.txt，写入 alpha 和 beta 两行内容
echo -e "alpha\nbeta" > "q01/input/docs/notes one.txt"

# 3. 创建隐藏文件 .secret.txt，写入内容 hidden
echo "hidden" > q01/input/docs/.secret.txt

# 4. 创建空文件 empty.txt
touch q01/input/tmp/empty.txt

# 5. 创建日志文件 run.log，写入两行自定义日志
echo -e "[INFO] task start\n[INFO] task finish" > q01/input/run.log

# 显示 q01 目录的绝对路径
realpath q01

# 长格式列出 input 目录下的全部项目，包括隐藏文件
ls -la q01/input/

mkdir -p work/25100021005

cd input

# 查找所有 .txt 文件，保留相对结构复制到目标目录
find . -name "*.txt" -exec cp --parents {} ../work/25100021005/ \;

# 批量设置所有子目录权限为 750
find work/25100021005 -type d -exec chmod 750 {} \;

# 批量设置所有普通文件权限为 640
find work/25100021005 -type f -exec chmod 640 {} \;

# 进入目标根目录，保证输出的是相对于该目录的相对路径
cd work/25100021005

# 查找所有文件，输出「相对路径 + 字节数」，写入 q01 目录下的 inventory.txt
find . -type f -printf "%P %s\n" > ../../inventory.txt
\end{lstlisting}

\begin{center}
{\large\bfseries 部分实验过程截图}
\end{center}

\resultimg{week1_check1.png}{题目一 实验过程截图}
\resultimg{week1_check1_2.png}{题目一 实验过程截图}

\begin{center}
{\large\bfseries 实验结果}
\end{center}
通过命令行在 q01 中完成了含空格文件名的批量整理：成功创建了所需目录与文件（含带空格的 notes one.txt 和隐藏文件），正确显示了绝对路径并列出全部文件（含隐藏文件），将 .txt 文件按原目录结构复制到 work/25100021005/，目录权限设为 750、文件设为 640，最终生成记录各文件相对路径与字节数的 inventory.txt。
\subsection*{题目二}
\resetpic
% 题目二 问题截图（代码汇总之前）
\resultimgcap[0.8\textwidth]{week1_q2.png}{题目二 题目要求}

\begin{lstlisting}[language=bash]
# TODO: 题目二代码汇总

# 一键生成 access.csv，写入题目给定的所有数据
# << 'EOF' ... EOF 是「原样写入」语法，不会转义特殊字符
cat > access.csv << 'EOF'
timestamp,user,path,status,latency_ms
09:00:01,alice,/api/users,200,120
09:00:02,bob,/api/orders,500,310
09:00:03,alice,/api/users,503,180
09:00:04,carol,/login,500,90
09:00:05,bob,/api/orders,502,260
09:00:06,dave,/health,200,20
09:00:07,alice,/api/orders,500,420
09:00:08,carol,/api/users,500,150
EOF

cat > analyze.sh << 'EOF'
#!/bin/bash

# ========== 1. 参数与文件存在性检查 ==========
# 如果没有传1个参数，输出用法到stderr，退出码1
if [ $# -ne 1 ]; then
    echo "用法: $0 <csv文件路径>" >&2
    exit 1
fi

csv_file="$1"

# 如果文件不存在，输出错误到stderr，退出码1（非零）
if [ ! -f "$csv_file" ]; then
    echo "错误: 文件 '$csv_file' 不存在" >&2
    exit 1
fi

# ========== 2. 统计 5xx 数量最多的前 2 个 path ==========
echo "5xx请求Top2路径:"
# 管道处理流程：
# 1. awk：逗号分隔，跳过表头，筛选状态码以5开头的行，输出第3列(path)
# 2. sort：排序，为后续去重统计做准备
# 3. uniq -c：统计每个路径出现的次数，格式为「次数 路径」
# 4. sort：先按第1列(次数)数字降序，次数相同按第2列(路径)字典序升序
# 5. head -n 2：取前2名
# 6. awk：只输出路径名
awk -F, 'NR>1 && $4 ~ /^5/ {print $3}' "$csv_file" \
    | sort \
    | uniq -c \
    | sort -k1,1nr -k2,2 \
    | head -n 2 \
    | awk '{print $2}'

# ========== 3. 计算平均延迟，保留两位小数 ==========
echo -e "\n平均延迟(ms):"
# NR>1 跳过表头；sum累加延迟，count计数；最后格式化输出平均值
awk -F, 'NR>1 {sum += $5; count++} END {printf "%.2f\n", sum / count}' "$csv_file"
EOF

# 给脚本添加可执行权限
chmod +x analyze.sh

#正常情况
./analyze.sh access.csv
#非正常情况（此时echo显示1）
./analyze.sh not_exist.csv
echo $?




\end{lstlisting}

\begin{center}
{\large\bfseries 部分实验过程截图}
\end{center}

\resultimg{week1_check2.png}{题目二 实验过程截图一}


\begin{center}
{\large\bfseries 实验结果}
\end{center}
编写 \texttt{analyze.sh} 脚本对 \texttt{access.csv} 进行日志统计：文件存在时正常处理，不存在时错误写入 stderr 并返回非零退出码。HTTP 5xx 数量最多的前 2 个 path 如图，按次数降序、字典序输出；全部数据行的平均 \texttt{latency\_ms} 如图。

\subsection*{题目三}
\resetpic
% 题目三 问题截图（代码汇总之前）
\resultimgcap[0.8\textwidth]{week1_q3.png}{题目三 题目要求}

\begin{lstlisting}[language=bash]
# TODO: 题目三代码汇总

#初始化仓库
git init

echo "mode=normal" > config.txt

#加入暂存区，并提交
git add config.txt
git commit -m "初始提交：创建config.txt，mode=normal"

#创建并切换到 `feature-a`
git checkout -b feature-a

#修改 config.txt 为 mode=safe，然后提交
echo "mode=safe" > config.txt
git add config.txtgit commit -m "feature-a：修改mode为safe"
#回到master，创建并切换到 `feature-b`，提交修改
git checkout master
git checkout -b feature-b
echo "mode=fast" > config.txt
git add config.txt
git commit -m "feature-b：修改mode为fast"

#回到master，先合并 feature-a
git checkout master
git merge feature-a -m "合并feature-a"

#合并b提示冲突，解决后提交
git merge feature-b
echo -e "mode=safe\nnote=reviewed" > config.txt
git add config.txt
git commit -m "解决合并冲突：保留mode=safe，新增note=reviewed"

#确认工作区干净
git status

#按题目要求查看完整分支历史图
git log --all --graph --decorate --oneline


\end{lstlisting}

\begin{center}
{\large\bfseries 部分实验过程截图}
\end{center}

\resultimg{week1_check3.png}{题目三 实验过程截图一}
\resultimg{week1_check3_2.png}{题目三 实验过程截图二（此图豆包提供，gitbash原截图没留存）}

\begin{center}
{\large\bfseries 实验结果}
\end{center}
在 q03 中从零创建了 Git 仓库并完成了合并冲突的制造与解决：main 分支上创建 config.txt（mode=normal）并完成初始提交；在 feature-a 分支改为 mode=safe 并提交，feature-b 分支改为 mode=fast 并提交。回到 main 先合并 feature-a、再合并 feature-b 时产生冲突，解决后保留 mode=safe 并新增一行 note=reviewed。最终工作区干净，git log --all --graph --decorate --oneline 显示了完整的分支与合并历史。

补充说明：本问题本身即是 Git 操作练习，其仓库仅用于制造并解决合并冲突，与最终提交到 GitHub 的报告仓库并不相同。为将实验报告提交到 GitHub，在合并冲突练习完成后，删除了该练习仓库的 \texttt{.git} 目录（避免将练习中的提交历史混入正式仓库），随后重新 \texttt{git init} 并关联远程仓库，再将报告相关内容提交上去。

\subsection*{题目四}
\resetpic
% 题目四 问题截图（代码汇总之前）
\resultimgcap[0.8\textwidth]{week1_q4.png}{题目四 题目要求}

\begin{lstlisting}[language=tex]
# TODO: 题目四代码汇总

\documentclass{article}
\usepackage{amsmath}
\usepackage{caption}
\usepackage[UTF8]{ctex}

\begin{document}
\title{Tool Report}
\author{曾奕豪 - 25100021005}
\maketitle

\section{Result}

% 正文引用公式：\eqref{标签名} 会自动生成带括号的公式编号
如式 \eqref{eq:sum} 所示，这是一个标准的自然数求和公式。
% equation 环境：自动生成公式编号
\begin{equation}
    \sum_{i=1}^{n} i = \frac{n(n+1)}{2}
    \label{eq:sum}  % \label 给公式打标签，用于交叉引用
\end{equation}

% 正文引用表格：\ref{标签名} 会自动生成表格编号
实验的核心参数统计如表 \ref{tab:params} 所示。
% table 浮动体环境：存放表格、标题和标签
\begin{table}[h]
    \centering % 表格居中
    \caption{实验核心参数表} % 表格标题
    \label{tab:params} % 表格标签，用于引用
    % 2行2列表格：{|c|c|} 表示两列居中、带竖线
    \begin{tabular}{|c|c|}
        \hline
        参数名称 & 参数数值 \\
        \hline
        平均延迟 & 193.75 ms \\
        \hline
    \end{tabular}
\end{table}
\end{document}

\end{lstlisting}

\begin{center}
{\large\bfseries 部分实验过程截图}
\end{center}

\resultimg{week1_check4.png}{题目四 实验过程截图一}
\resultimg{week1_check4_2.png}{题目四 实验过程截图二}

\begin{center}
{\large\bfseries 实验结果}
\end{center}
将模板保存为 q04/report.tex 并补全内容后，用 latexmk -pdf -halt-on-error 成功构建出 report.pdf。文档中加入了带编号和 label 的公式，并在正文通过 ref/eqref 命令正确引用；加入了一个 2 行 2 列的表格并设置 caption 和 label，正文通过 ref 命令正确引用该表。生成的 PDF 可正常打开，交叉引用均正确显示，未出现 "??" 占位符。

\subsection*{课上阶段提交记录}
% 课上 4 题的 commit 记录截图（git log 或仓库 commits 列表）
\resultimgcap[0.9\textwidth]{week1_commit_kecheng.png}{课上阶段 commit 记录}

% ---------- 2. 课后练习与结果 ----------
\section{课后练习与结果}
% 说明：课后完成 >10 个实例，每个实例按"题目 → 代码 → 运行结果"组织。

\subsection*{实例 1}
\resetpic
\begin{lstlisting}
# TODO: 实例1代码汇总
\end{lstlisting}

\begin{center}
{\large\bfseries 部分实验过程截图}
\end{center}

\resultimg{week1_ex1.png}{实例一 运行结果图一}
\resultimg{week1_ex1_2.png}{实例一 运行结果图二}

\begin{center}
{\large\bfseries 实验结果}
\end{center}

\subsection*{实例 2}
\resetpic
% TODO: 继续补充实例 2 ~ 实例 10+，结构与实例 1 相同
\begin{lstlisting}
# TODO: 实例2代码汇总
\end{lstlisting}

\begin{center}
{\large\bfseries 部分实验过程截图}
\end{center}

\resultimg{week1_ex2.png}{实例二 运行结果图一}
\resultimg{week1_ex2_2.png}{实例二 运行结果图二}

\begin{center}
{\large\bfseries 实验结果}
\end{center}

\subsection*{课后阶段提交记录}
% 课后练习的 commit 记录截图
\resultimgcap[0.9\textwidth]{week1_commit_kehou.png}{课后阶段 commit 记录}

% ---------- 3. 解题感悟 ----------
\section{解题感悟}
% TODO: 记录做题过程中的思考、遇到的问题与解决方式、收获与总结。
此处填写解题感悟。



\end{document}
